\scanpage{24}
\diagram{\begin{tikzcd}[column sep=large]
\operatorname{dét}T_p(A)_{\text{ét}}\arrow[r,"\sim"] & \operatorname{dét}T_p(A^*)
\end{tikzcd}}
[\deleted{Soit} Si on a $\ell\ne p$ au lieu de $p$, on a bien
\diagram{\begin{tikzcd}[column sep=large]
\operatorname{dét}T_\ell(A)\arrow[r,phantom,"\simeq" description,"\alpha_\ell"] & \operatorname{dét}T_\ell(A^*)\arrow[r,phantom,"\simeq" description] & \mathbb Z_\ell(g)
\end{tikzcd}}]
\deleted{car $T_\ell(A^*)\simeq\deleted{\operatorname{dét}}T_\ell(\unsure{A^*|A})^\vee\otimes\mathbb Z_\ell(1)$ donc $\operatorname{dét}T_\ell(A^*)\simeq\operatorname{dét}T_\ell(A)$ (Choisissons une polarisation)}

Si $\varphi$ est une polarisation de $A$, définissant
\diagram{\begin{tikzcd}\widetilde\varphi:A\arrow[r,dash] & A^*\end{tikzcd}}
d'où
\diagram{\begin{tikzcd}T_\ell(\widetilde\varphi):T_\ell(A)\arrow[r,dash] & T_\ell(A^*)\end{tikzcd}},
\diagram{\begin{tikzcd}[column sep=large]
\operatorname{dét}T_\ell(\widetilde\varphi):\operatorname{dét}T_\ell(A)\arrow[r,"\beta_\ell(\varphi)"] & \operatorname{dét}T_\ell(A^*)
\end{tikzcd}},
alors l'isom $\alpha_\ell$ n'est autre que $\beta_\ell(\varphi)/\deg\varphi$. En s'inspirant de cette formule, on définira;
\diagram{\begin{tikzcd}[column sep=large]
\beta_p(\varphi)=\operatorname{dét}(T_p(\widetilde\varphi)_{\text{ét}}):\operatorname{dét}T_p(A)_{\text{ét}}\arrow[r] & \operatorname{dét}T_p(A^*)_{\text{ét}}
\end{tikzcd}}
et on considère
\[
\alpha_p(\varphi)=\beta_p(\varphi)/\sqrt{\deg\widetilde\varphi}:
\]
\diagram{\begin{tikzcd}[column sep=large]
\operatorname{dét}T_p(A)_{\text{ét}}\otimes\mathbb Q_{\unsure{\ell|p}} & \operatorname{dét}T_p(A^*)_{\text{ét}}\inserted{\otimes\mathbb Q_{\unsure{\ell|p}}}\arrow[l]
\end{tikzcd}}
À prouver: a) $\alpha_p(\varphi)$ ne dépend pas de $\varphi$.

\deleted{est \unsure{pratiquement} formelle de ceci}
\deleted{\uline{Lemme} Soit $u$ un \deleted{\unsure{autom}}\deleted{\unsure{endom}}\inserted{endom} de $A$, d'où
\diagram{\begin{tikzcd}[column sep=large]
[\operatorname{dét}T_p(u)_{\text{ét}}:\operatorname{dét}T_p(A)_{\text{ét}}\deleted{\unsure{\mathbf Z_p}}\arrow[r] & \operatorname{dét}T_p(A)_{\text{ét}}]\in\mathbb Z_p.
\end{tikzcd}}
On a
\[
\left\{\begin{aligned}
&\deg u\text{ est un carré parfait}\\
&\operatorname{dét}(T_p(u)_{\text{ét}})=\sqrt{\deg u}
\end{aligned}\right.
\]
i.e.
\[
\deg u=(\operatorname{dét}(T_p(u)_{\text{ét}}))^2
\]
\uline{et} $\operatorname{dét}T_p(u)_{\text{ét}}$ est un entier $>0$}
\margin{\uline{faux}}
\mspage{23}
